\documentclass[a4paper,12pt]{article}

% ============================================================
% КОДИРОВКА И ЯЗЫК
% ============================================================
\usepackage[T2A]{fontenc}
\usepackage[utf8]{inputenc}
\usepackage[russian]{babel}

% ============================================================
% МАТЕМАТИКА И ТИПОГРАФИКА
% ============================================================
\usepackage{amsmath,amssymb,amsthm,mathtools}
\usepackage{helvet}
\renewcommand{\familydefault}{\sfdefault}
\usepackage{icomma}
\usepackage[a4paper,left=20mm,right=20mm,top=20mm,bottom=22mm]{geometry}
\usepackage{microtype}
\usepackage{parskip}
\usepackage{setspace}
\usepackage{xcolor}

% ============================================================
% ГРАФИКА
% ============================================================
\usepackage{tikz}
\usetikzlibrary{calc}
\usetikzlibrary{arrows.meta,positioning,shapes.geometric}

\usepackage{tkz-euclide}

\usepackage{pgfplots}
\pgfplotsset{compat=1.18}

% ============================================================
% ТАБЛИЦЫ И СПИСКИ
% ============================================================
\usepackage{tabularx,booktabs,longtable}
\usepackage{enumitem}

% ============================================================
% СЛУЖЕБНЫЕ ПАКЕТЫ
% ============================================================
\usepackage{hyperref}
\usepackage{fancyhdr}

% ============================================================
% ЦВЕТОВАЯ ПАЛИТРА
% ============================================================
\definecolor{accent}{HTML}{008080}
\definecolor{darkaccent}{HTML}{004D4D}
\definecolor{textgray}{HTML}{333333}
\definecolor{softgray}{HTML}{777777}
\definecolor{hiddenedge}{HTML}{9A9A9A}

\hypersetup{
    colorlinks=true,
    linkcolor=darkaccent,
    urlcolor=darkaccent
}

\color{textgray}
\onehalfspacing

\setlength{\parindent}{0pt}
\setlength{\parskip}{0.55em}

\setlist[itemize]{
    leftmargin=1.6em,
    itemsep=0.25em,
    topsep=0.3em
}

\setlist[enumerate]{
    leftmargin=1.8em,
    itemsep=0.45em,
    topsep=0.4em
}

% ============================================================
% КОЛОНТИТУЛЫ
% ============================================================
\pagestyle{fancy}
\fancyhf{}

\fancyhead[L]{\small\color{softgray}Кристина}
\fancyhead[R]{\small\color{softgray}ЕГЭ. Профильная математика}
\fancyfoot[C]{\small\color{softgray}\thepage}

\renewcommand{\headrulewidth}{0.4pt}
\renewcommand{\headrule}{%
    \hbox to\headwidth{%
        \color{accent}\leaders\hrule height \headrulewidth\hfill
    }%
}

\setlength{\headheight}{15pt}

% ============================================================
% ЗАГОЛОВКИ
% ============================================================
\newcommand{\sectiontitle}[1]{%
    \vspace{0.6em}
    {\Large\bfseries\color{darkaccent}#1}\par
    \vspace{0.25em}
    {\color{accent}\rule{\linewidth}{0.7pt}}
    \vspace{0.45em}
}

\newcommand{\subsectiontitle}[1]{%
    \vspace{0.4em}
    {\large\bfseries\color{darkaccent}#1}\par
    \vspace{0.2em}
}

\newcommand{\keyidea}[1]{%
    \par\medskip
    {\bfseries\color{darkaccent}Ключевая идея.} #1
    \par\medskip
}

% ============================================================
% TKZ-EUCLIDE
% ============================================================
\tkzSetUpLine[
    line width=1pt,
    color=textgray
]

\tkzSetUpPoint[
    color=darkaccent,
    fill=darkaccent,
    size=2.8
]

\tkzSetUpLabel[
    color=black,
    font=\small
]

% ============================================================
% СТИЛИ БЛОК-СХЕМ
% ============================================================
\tikzset{
    flowstart/.style={
        ellipse,
        draw=darkaccent,
        line width=0.9pt,
        align=center,
        minimum width=36mm,
        minimum height=11mm,
        inner sep=4pt
    },
    flowaction/.style={
        rounded corners=3pt,
        rectangle,
        draw=accent,
        line width=0.9pt,
        align=center,
        text width=92mm,
        minimum height=13mm,
        inner sep=5pt
    },
    flowdecision/.style={
        diamond,
        aspect=2.2,
        draw=darkaccent,
        line width=0.9pt,
        align=center,
        text width=40mm,
        inner sep=2pt
    },
    flowarrow/.style={
        -{Stealth[length=3mm,width=2mm]},
        line width=0.9pt,
        color=darkaccent
    }
}

\begin{document}

% ============================================================
% ТИТУЛЬНЫЙ ЛИСТ
% ============================================================
\begin{titlepage}
\thispagestyle{empty}

\begin{center}

\vspace*{23mm}

{\small\color{softgray}
ПРОФИЛЬНАЯ МАТЕМАТИКА
\textbullet\
10 КЛАСС
\textbullet\
ЕГЭ
}

\vspace{15mm}

{\Huge\bfseries\color{darkaccent}
Чек-лист
}

\vspace{7mm}

{\LARGE\bfseries
Стереометрия: куб\\[2mm]
формулы, диагонали и моделирование задач
}

\vspace{15mm}

{\large
Ключевые формулы куба,\\
главная диагональ и диагональ грани,\\
перевод условия на язык уравнений
}

\vfill

{\large\bfseries
Кристина
}

\vspace{7mm}

{\normalsize
Лёвин Артём Александрович\\
эксклюзивно для Кристины
}

\vspace{5mm}

{\normalsize
\today
}

\vspace{24mm}

\begin{tikzpicture}[x=1cm,y=1cm]
    \draw[
        accent,
        line width=1.15pt
    ]
    (0,0)
    .. controls (3,0.8) and (7,-0.8) ..
    (10,0)
    .. controls (12,0.45) and (14,0.35) ..
    (16,0.05);
\end{tikzpicture}

\end{center}
\end{titlepage}

% ============================================================
% 1. ОСНОВНЫЕ ФОРМУЛЫ
% ============================================================
\sectiontitle{1. Что нужно знать о кубе}

Куб --- прямоугольный параллелепипед, у которого все рёбра равны.
Обозначим длину ребра через \(a\).

Для геометрической фигуры обязательно
\[
a>0.
\]

Основные величины:

\begin{center}
\renewcommand{\arraystretch}{1.4}

\begin{tabularx}{0.94\linewidth}{@{}>{\bfseries}l X c@{}}
\toprule

Величина
&
Смысл
&
Формула
\\

\midrule

Объём
&
пространство внутри куба
&
\(V=a^3\)
\\

Площадь полной поверхности
&
сумма площадей шести квадратных граней
&
\(S_{\text{полн}}=6a^2\)
\\

Диагональ грани
&
диагональ одной квадратной грани
&
\(d_{\text{гр}}=a\sqrt2\)
\\

Диагональ куба
&
отрезок между противоположными вершинами куба
&
\(d=a\sqrt3\)
\\

\bottomrule
\end{tabularx}
\end{center}

\keyidea{
Если в условии сказано просто «диагональ куба», используется
пространственная диагональ \(d=a\sqrt3\).
Фраза «диагональ грани» означает диагональ квадратной грани
\(d_{\text{гр}}=a\sqrt2\).
}

% ============================================================
% ПОЧЕМУ КОРНИ 2 И 3
% ============================================================
\subsectiontitle{Почему появляются \(\sqrt2\) и \(\sqrt3\)}

На квадратной грани по теореме Пифагора:
\[
d_{\text{гр}}^2=a^2+a^2=2a^2.
\]

Так как длина положительна,
\[
d_{\text{гр}}=a\sqrt2.
\]

Для пространственной диагонали:
\[
d^2=d_{\text{гр}}^2+a^2.
\]

Следовательно,
\[
d^2=2a^2+a^2=3a^2,
\]
поэтому
\[
d=a\sqrt3.
\]

% ============================================================
% ЧЕРТЁЖ КУБА
% ============================================================
\subsectiontitle{Схема куба}

\begin{center}

\begin{tikzpicture}[scale=1.08]

% ------------------------------------------------------------
% Вершины
% Передняя грань: ABCD
% Задняя грань: A_1B_1C_1D_1
% ------------------------------------------------------------

\tkzDefPoint(0,0){A}
\tkzDefPoint(5.2,0){B}
\tkzDefPoint(5.2,4.7){C}
\tkzDefPoint(0,4.7){D}

\tkzDefPoint(1.65,1.35){A1}
\tkzDefPoint(6.85,1.35){B1}
\tkzDefPoint(6.85,6.05){C1}
\tkzDefPoint(1.65,6.05){D1}

% ------------------------------------------------------------
% Видимые рёбра
% ------------------------------------------------------------

\tkzDrawSegments[
    color=textgray,
    line width=1.05pt
](A,B B,C C,D D,A)

\tkzDrawSegments[
    color=textgray,
    line width=1.05pt
](B,B1 B1,C1 C1,D1 D1,D C,C1)

% ------------------------------------------------------------
% Скрытые рёбра
% ------------------------------------------------------------

\tkzDrawSegments[
    color=hiddenedge,
    line width=0.9pt,
    dashed
](A,A1 A1,B1 A1,D1)

% ------------------------------------------------------------
% Диагональ грани A_1B
% ------------------------------------------------------------

\tkzDrawSegment[
    color=darkaccent,
    line width=1.8pt
](A1,B)

% ------------------------------------------------------------
% Пространственная диагональ A_1C
% ------------------------------------------------------------

\tkzDrawSegment[
    color=accent,
    line width=2.0pt
](A1,C)

% ------------------------------------------------------------
% Точки
% ------------------------------------------------------------

\tkzDrawPoints(A,B,C,D,A1,B1,C1,D1)

% ------------------------------------------------------------
% Подписи вершин
% ------------------------------------------------------------

\tkzLabelPoint[below left](A){$A$}
\tkzLabelPoint[below right](B){$B$}
\tkzLabelPoint[right](C){$C$}
\tkzLabelPoint[left](D){$D$}

\tkzLabelPoint[above left](A1){$A_1$}
\tkzLabelPoint[right](B1){$B_1$}
\tkzLabelPoint[above right](C1){$C_1$}
\tkzLabelPoint[above left](D1){$D_1$}

% ------------------------------------------------------------
% Подпись ребра
% ------------------------------------------------------------

\node[
    below=4pt,
    color=textgray,
    font=\small
]
at ($(A)!0.5!(B)$)
{$a$};

% ------------------------------------------------------------
% Подпись диагонали грани A_1B
% ------------------------------------------------------------

\node[
    color=darkaccent,
    fill=white,
    inner sep=2.5pt,
    font=\small
]
at ($(A1)!0.47!(B)+(0,-0.42)$)
{$A_1B=d_{\text{гр}}=a\sqrt2$};

% ------------------------------------------------------------
% Подпись пространственной диагонали A_1C
% ------------------------------------------------------------

\node[
    color=accent,
    fill=white,
    inner sep=2.5pt,
    font=\small
]
at ($(A1)!0.58!(C)+(0.30,0.38)$)
{$A_1C=d=a\sqrt3$};

\end{tikzpicture}

\end{center}

\keyidea{
На чертеже \(A_1B\) --- диагональ нижней грани \(AA_1B_1B\),
а \(A_1C\) --- пространственная диагональ куба.
}

% ============================================================
% 2. БАЗОВЫЙ МЕТОД
% ============================================================
\sectiontitle{2. Базовый метод: начать с искомой величины}

Если задача просит найти величину, для которой известна формула,
удобно сразу записать эту формулу.

Например, если требуется найти объём:
\[
\boxed{V=a^3}.
\]

После этого определяем, какой величины не хватает.

Если неизвестно ребро \(a\), ищем формулу,
которая связывает \(a\) с данными задачи.

% ============================================================
% ПРИМЕР 1
% ============================================================
\subsectiontitle{Пример 1. Известна диагональ куба}

Диагональ куба равна
\[
d=\sqrt{27}.
\]

Найдите объём куба.

Начинаем с искомой величины:
\[
V=a^3.
\]

Для вычисления объёма нужно найти \(a\).

Используем
\[
d=a\sqrt3.
\]

Подставляем:
\[
\sqrt{27}=a\sqrt3.
\]

Преобразуем корень:
\[
\sqrt{27}
=
\sqrt{9\cdot3}
=
3\sqrt3.
\]

Получаем
\[
3\sqrt3=a\sqrt3.
\]

Так как
\[
\sqrt3\ne0,
\]
можно разделить обе части на \(\sqrt3\):
\[
a=3.
\]

Тогда
\[
V=3^3=27.
\]

\[
\boxed{V=27}
\]

\subsectiontitle{Проверка ограничения}

Ребро куба должно удовлетворять условию
\[
a>0.
\]

Полученное значение
\[
a=3
\]
условию удовлетворяет.

% ============================================================
% БЛОК-СХЕМА 1
% ============================================================
\clearpage
\thispagestyle{fancy}

\begin{center}

{\LARGE\bfseries\color{darkaccent}
Алгоритм: задача на формульную величину куба
}

\end{center}

\vspace{10mm}

\begin{center}

\begin{tikzpicture}[node distance=8mm]

\node[flowstart] (start)
{Начало};

\node[flowaction,below=of start] (target)
{Запишите формулу величины, которую требуется найти};

\node[flowdecision,below=11mm of target] (enough)
{Все необходимые\\величины известны?};

\node[flowaction,below=12mm of enough] (find)
{Выберите формулу, связывающую неизвестную величину с данными задачи};

\node[flowaction,below=of find] (solve)
{Подставьте данные и найдите недостающую величину};

\node[flowaction,below=of solve] (restrict)
{Проверьте геометрические ограничения};

\node[flowaction,below=of restrict] (substitute)
{Подставьте найденное значение в исходную формулу};

\node[flowstart,below=of substitute] (finish)
{Ответ};

\draw[flowarrow]
(start) -- (target);

\draw[flowarrow]
(target) -- (enough);

\draw[flowarrow]
(enough) --
node[right,font=\small]{нет}
(find);

\draw[flowarrow]
(find) -- (solve);

\draw[flowarrow]
(solve) -- (restrict);

\draw[flowarrow]
(restrict) -- (substitute);

\draw[flowarrow]
(substitute) -- (finish);

\draw[flowarrow]
(enough.east)
-- ++(20mm,0)
|- node[pos=0.22,right,font=\small]{да}
(substitute.east);

\end{tikzpicture}

\end{center}

\clearpage

% ============================================================
% 3. ДВА ВИДА ДИАГОНАЛЕЙ
% ============================================================
\sectiontitle{3. Диагональ грани и диагональ куба}

\subsectiontitle{Пример 2. Известна диагональ грани}

Диагональ грани куба равна
\[
d_{\text{гр}}=\sqrt8.
\]

Найдите объём.

Искомая величина:
\[
V=a^3.
\]

Используем формулу именно для диагонали грани:
\[
d_{\text{гр}}=a\sqrt2.
\]

Подставляем:
\[
\sqrt8=a\sqrt2.
\]

Поскольку
\[
\sqrt8
=
\sqrt{4\cdot2}
=
2\sqrt2,
\]
получаем
\[
2\sqrt2=a\sqrt2.
\]

Так как
\[
\sqrt2\ne0,
\]
получаем
\[
a=2.
\]

Следовательно,
\[
V=2^3=8.
\]

\[
\boxed{V=8}
\]

% ============================================================
% СРАВНЕНИЕ
% ============================================================
\subsectiontitle{Сравнение двух ситуаций}

\begin{center}

\renewcommand{\arraystretch}{1.45}

\begin{tabularx}{\linewidth}{@{}X X@{}}
\toprule

\textbf{Диагональ грани}
&
\textbf{Диагональ куба}
\\

\midrule

\[
d_{\text{гр}}=a\sqrt2
\]
&
\[
d=a\sqrt3
\]
\\

лежит в одной квадратной грани
&
соединяет противоположные вершины куба
\\

для поиска ребра:
\[
a=\frac{d_{\text{гр}}}{\sqrt2}
\]
&
для поиска ребра:
\[
a=\frac{d}{\sqrt3}
\]
\\

\bottomrule
\end{tabularx}

\end{center}

\keyidea{
Перед вычислениями определите вид диагонали.
Подмена \(\sqrt2\) на \(\sqrt3\) или наоборот полностью меняет решение.
}

% ============================================================
% 4. БЫЛО — СТАЛО
% ============================================================
\sectiontitle{4. Модель «было --- стало»}

В более сложных задачах данные часто задаются словами:

\begin{itemize}
    \item ребро увеличили на некоторое число;
    \item объём увеличился на некоторое число;
    \item площадь уменьшилась;
    \item длина стала в несколько раз больше.
\end{itemize}

Такие формулировки нужно перевести на математический язык.

Пусть
\[
a_1
\]
--- первоначальное ребро,

\[
a_2
\]
--- новое ребро,

\[
V_1
\]
--- первоначальный объём,

\[
V_2
\]
--- новый объём.

Если ребро увеличили на \(1\), то
\[
\boxed{a_2=a_1+1}.
\]

Если объём увеличился на \(19\), то
\[
\boxed{V_2=V_1+19}.
\]

Поскольку
\[
V_1=a_1^3,
\qquad
V_2=a_2^3,
\]
получаем
\[
\begin{cases}
a_2=a_1+1,\\[2mm]
a_2^3=a_1^3+19.
\end{cases}
\]

\keyidea{
Фраза «увеличилось на \(k\)» означает
\[
\text{новое значение}
=
\text{старое значение}+k.
\]
}

% ============================================================
% ПРИМЕР 3
% ============================================================
\subsectiontitle{Пример 3. Ребро увеличили на 1}

Если каждое ребро куба увеличить на \(1\),
его объём увеличится на \(19\).

Найдите первоначальное ребро.

Записываем:
\[
a_2=a_1+1,
\]
\[
V_2=V_1+19.
\]

Используем формулу объёма:
\[
a_2^3=a_1^3+19.
\]

Подставляем
\[
a_2=a_1+1:
\]
\[
(a_1+1)^3=a_1^3+19.
\]

Используем формулу куба суммы:
\[
(x+y)^3
=
x^3+3x^2y+3xy^2+y^3.
\]

При \(y=1\):
\[
(a_1+1)^3
=
a_1^3+3a_1^2+3a_1+1.
\]

Получаем:
\[
a_1^3+3a_1^2+3a_1+1
=
a_1^3+19.
\]

Сокращаем одинаковые слагаемые:
\[
3a_1^2+3a_1+1=19.
\]

Переносим:
\[
3a_1^2+3a_1-18=0.
\]

Делим на \(3\):
\[
a_1^2+a_1-6=0.
\]

Разложим на множители:
\[
(a_1+3)(a_1-2)=0.
\]

Получаем:
\[
a_1=-3
\qquad\text{или}\qquad
a_1=2.
\]

По геометрическому смыслу
\[
a_1>0.
\]

Поэтому
\[
\boxed{a_1=2}.
\]

% ============================================================
% БЛОК-СХЕМА 2
% ============================================================
\clearpage
\thispagestyle{fancy}

\begin{center}

{\LARGE\bfseries\color{darkaccent}
Алгоритм: задачи «было --- стало»
}

\end{center}

\vspace{8mm}

\begin{center}

\begin{tikzpicture}[node distance=7mm]

\node[flowstart] (start)
{Начало};

\node[flowaction,below=of start] (names)
{Введите обозначения для первоначальных и новых величин};

\node[flowaction,below=of names] (translate)
{Переведите слова условия в математические равенства};

\node[flowaction,below=of translate] (formulas)
{Замените геометрические величины их формулами};

\node[flowdecision,below=11mm of formulas] (onevar)
{Получено уравнение\\с одной неизвестной?};

\node[flowaction,below=12mm of onevar] (substitution)
{Используйте дополнительную связь из условия и выполните подстановку};

\node[flowaction,below=of substitution] (solve)
{Решите полученное уравнение};

\node[flowaction,below=of solve] (check)
{Отбросьте значения, противоречащие геометрическому смыслу};

\node[flowstart,below=of check] (finish)
{Ответ};

\draw[flowarrow]
(start) -- (names);

\draw[flowarrow]
(names) -- (translate);

\draw[flowarrow]
(translate) -- (formulas);

\draw[flowarrow]
(formulas) -- (onevar);

\draw[flowarrow]
(onevar) --
node[right,font=\small]{нет}
(substitution);

\draw[flowarrow]
(substitution) -- (solve);

\draw[flowarrow]
(solve) -- (check);

\draw[flowarrow]
(check) -- (finish);

\draw[flowarrow]
(onevar.east)
-- ++(21mm,0)
|- node[pos=0.22,right,font=\small]{да}
(solve.east);

\end{tikzpicture}

\end{center}

\clearpage

% ============================================================
% 5. КУБ СУММЫ
% ============================================================
\sectiontitle{5. Куб суммы и полезная мнемоника}

В задачах типа «ребро увеличили на \(1\)» естественно возникает выражение
\[
(a+1)^3.
\]

Формула:
\[
\boxed{
(x+y)^3
=
x^3+3x^2y+3xy^2+y^3
}
\]

Коэффициенты идут в порядке
\[
1,\;3,\;3,\;1.
\]

Для быстрого воспроизведения удобно запомнить:
\[
\boxed{1\;3\;3\;1}.
\]

При \(y=1\):
\[
(a+1)^3
=
a^3+3a^2+3a+1.
\]

Для разности:
\[
\boxed{
(x-y)^3
=
x^3-3x^2y+3xy^2-y^3
}
\]

Знаки чередуются:
\[
+,\;-,\;+,\;-.
\]

\subsectiontitle{Типичная ошибка}

Неверно:
\[
(a+1)^3=a^3+1.
\]

Правильно:
\[
(a+1)^3
=
a^3+3a^2+3a+1.
\]

% ============================================================
% 6. ПЛОЩАДЬ ЧЕРЕЗ ДИАГОНАЛЬ
% ============================================================
\sectiontitle{6. Площадь поверхности через диагональ}

\subsectiontitle{Пример 4}

Диагональ куба равна \(1\).

Найдите площадь полной поверхности.

Начинаем с формулы искомой величины:
\[
S_{\text{полн}}=6a^2.
\]

Ребро неизвестно.

Используем:
\[
d=a\sqrt3.
\]

Так как
\[
d=1,
\]
получаем
\[
1=a\sqrt3.
\]

Следовательно,
\[
a=\frac{1}{\sqrt3}.
\]

Поскольку
\[
\sqrt3>0,
\]
деление допустимо.

Подставляем:
\[
S_{\text{полн}}
=
6\left(\frac{1}{\sqrt3}\right)^2.
\]

Получаем:
\[
S_{\text{полн}}
=
6\cdot\frac13
=
2.
\]

\[
\boxed{S_{\text{полн}}=2}
\]

% ============================================================
% 7. ТИПИЧНЫЕ ОШИБКИ
% ============================================================
\sectiontitle{7. Что проверить перед ответом}

\begin{enumerate}

\item
\textbf{Какая именно диагональ дана?}

\[
d_{\text{гр}}=a\sqrt2,
\qquad
d=a\sqrt3.
\]

\item
\textbf{Какую величину требуется найти?}

Если объём:
\[
V=a^3.
\]

Если площадь полной поверхности:
\[
S_{\text{полн}}=6a^2.
\]

\item
\textbf{Переведены ли слова условия в математические равенства?}

Например:
\[
\text{«увеличили на }1\text{»}
\quad\Longrightarrow\quad
a_2=a_1+1.
\]

\item
\textbf{Правильно ли раскрыт куб суммы?}

\[
(x+y)^3
=
x^3+3x^2y+3xy^2+y^3.
\]

\item
\textbf{Проверен ли геометрический смысл корней?}

Для ребра:
\[
a>0.
\]

Для диагоналей:
\[
d>0,
\qquad
d_{\text{гр}}>0.
\]

\item
\textbf{Правильно ли преобразованы корни?}

Например:
\[
\sqrt{27}=3\sqrt3,
\qquad
\sqrt8=2\sqrt2.
\]

\end{enumerate}

% ============================================================
% 8. ИТОГ
% ============================================================
\sectiontitle{8. Итоговый чек-лист}

Перед решением задачи о кубе полезно быстро воспроизвести четыре формулы:

\[
\boxed{V=a^3}
\]

\[
\boxed{S_{\text{полн}}=6a^2}
\]

\[
\boxed{d_{\text{гр}}=a\sqrt2}
\]

\[
\boxed{d=a\sqrt3}
\]

Также нужна формула:
\[
\boxed{
(x+y)^3
=
x^3+3x^2y+3xy^2+y^3
}
\]

Практический порядок решения:

\begin{enumerate}

\item
внимательно определить, что дано и что требуется найти;

\item
записать формулу искомой величины;

\item
определить недостающую величину;

\item
выбрать формулу, связывающую её с данными;

\item
выполнить преобразования;

\item
проверить ограничения;

\item
записать ответ.

\end{enumerate}

\keyidea{
В составной задаче полезно стремиться к уравнению с одной неизвестной.
}

% ============================================================
% ТРЕНИРОВОЧНЫЙ БЛОК
% ============================================================
\clearpage
\thispagestyle{fancy}

\begin{center}

{\LARGE\bfseries\color{darkaccent}
Тренировочный блок
}

\vspace{2mm}

{\large
Путь А
\textbullet\
5 заданий
}

\end{center}

\vspace{7mm}

\textbf{1.}
Диагональ куба равна
\[
4\sqrt3.
\]
Найдите объём куба.

\vspace{5mm}

\textbf{2.}
Диагональ грани куба равна
\[
6\sqrt2.
\]
Найдите площадь полной поверхности куба.

\vspace{5mm}

\textbf{3.}
Объём куба равен
\[
125.
\]
Найдите диагональ куба.

\vspace{5mm}

\textbf{4.}
Диагональ куба равна
\[
3.
\]
Найдите площадь полной поверхности куба.

\vspace{5mm}

\textbf{5.}
Если каждое ребро куба увеличить на \(1\),
объём куба увеличится на \(37\).

Найдите первоначальную длину ребра.

\vfill

{\small\color{softgray}
При выполнении заданий начинайте с формулы искомой величины.
В задаче 5 используйте модель «было --- стало».
}

% ============================================================
% ОТВЕТЫ
% ============================================================
\clearpage
\thispagestyle{fancy}

\begin{center}

{\LARGE\bfseries\color{darkaccent}
Ответы к тренировочному блоку
}

\end{center}

\vspace{9mm}

\begin{table}[h]
\centering
\caption{Ответы}

\renewcommand{\arraystretch}{1.45}

\begin{tabularx}{0.88\linewidth}{@{}c X@{}}
\toprule

\textbf{№}
&
\textbf{Ответ}
\\

\midrule

1
&
\[
V=64
\]
\\

2
&
\[
S_{\text{полн}}=216
\]
\\

3
&
\[
d=5\sqrt3
\]
\\

4
&
\[
S_{\text{полн}}=18
\]
\\

5
&
\[
a=3
\]
\\

\bottomrule
\end{tabularx}

\end{table}

\vfill

\begin{center}

{\small\color{softgray}
Лёвин Артём Александрович
\textbullet\
профильная математика
}

\end{center}

\end{document}